\begin{table}[ht]
\centering
\caption{Multipliers under Alternative Fiscal Rules and the Capital-Goods Friction}
\label{tab:mult_compare}
\begin{tabular}{lcccccccc}
\hline\hline
 & \multicolumn{4}{c}{$H=12$} & \multicolumn{4}{c}{$H=20$} \\
\cmidrule(lr){2-5}\cmidrule(lr){6-9}
Instrument & FB & AR(1) & FB+K & AR(1)+K & FB & AR(1) & FB+K & AR(1)+K \\
\hline
Government spending      & 0.42 & 0.70 & 0.50 & 0.74 & $-$0.26 & 0.57 & $-$0.09 & 0.64 \\
Labor tax                & 0.65 & 0.72 & 0.55 & 0.67 & 0.18 & 0.83 & 0.15 & 0.76 \\
Capital tax              & 2.82 & 2.82 & 2.69 & 2.67 & 2.70 & 3.02 & 2.65 & 2.81 \\
Consumption tax          & 0.03 & 0.34 & $-$0.01 & 0.33 & $-$0.62 & 0.33 & $-$0.57 & 0.34 \\
Investment tax credit    & 2.23 & 1.43 & 1.68 & 1.19 & 2.14 & 1.77 & 1.51 & 1.36 \\
Expensing rate           & 0.19 & 0.32 & 0.07 & 0.31 & $-$0.37 & 0.44 & $-$0.48 & 0.38 \\
\hline\hline
\multicolumn{9}{p{0.95\textwidth}}{\footnotesize Note: Posterior medians of the present-value own-revenue multiplier. FB: baseline fiscal rules with debt and output feedback. AR(1): only transfers respond to debt. +K: capital-goods friction with $\epsilon_K = 4$. Each column is a separately estimated model. Intervals are in Tables~\ref{tab:mult_main} and~\ref{tab:mult_alt}.}
\end{tabular}
\end{table}
